\section{Selective Prediction and Rationale Faithfulness}
\label{sec:selective-salience}

\subsection{Selective prediction objective}
FactKG labels are binary (SUPPORTED/REFUTED), but practical deployment requires controlled abstention when evidence sufficiency is low.
M3 computes an abstain score per claim and supports both:
\begin{itemize}
\item threshold-free risk-coverage behavior (\AURC{}),
\item optional fixed operating point via \texttt{tau\_abstain}.
\end{itemize}
This follows the selective classification framing where a model may abstain to improve reliability at chosen coverage levels \cite{geifman2017selective}.

\subsection{Abstain score used in code}
For class probabilities $\hat{\mathbf{p}}$ and sufficiency score $\texttt{esi\_geom}$:
\[
\texttt{abstain\_score}
= \alpha\cdot (1-\max(\hat{\mathbf{p}})) + (1-\alpha)\cdot (1-\texttt{esi\_geom}),
\]
with default $\alpha=0.5$.

\subsection{Risk-coverage evaluation and \AURC}
Given abstain scores and correctness labels:
\begin{enumerate}
\item sort claims by ascending abstain score (most answerable first),
\item reveal predictions progressively,
\item at each coverage level compute risk $=$ error rate among answered claims.
\end{enumerate}
Area under risk-coverage curve is computed by trapezoidal integration:
\[
\AURC = \sum_i (\phi_i-\phi_{i-1})\cdot\frac{R_i+R_{i-1}}{2}.
\]
Lower \AURC{} is better.

\subsection{Fixed operating point (\texttt{tau\_abstain})}
When \texttt{tau\_abstain} is set, claims with
\[
\texttt{abstain\_score} \le \texttt{tau\_abstain}
\]
are answered, and the rest are abstained.
Logged metrics:
\begin{itemize}
\item \texttt{coverage} $=$ answered / total,
\item \texttt{risk} $=$ errors among answered / answered,
\item \texttt{accuracy\_on\_answered},
\item \texttt{abstain\_rate} $=1-\texttt{coverage}$.
\end{itemize}

\subsection{Rationale edge extraction}
\texttt{extract\_top\_rationale\_edges} ranks edges by:
\[
\text{salience}_{\text{attn-pv}}(e)=\texttt{attention\_used}(e)\cdot \Rel(e).
\]
where \texttt{attention\_used} is class-conditioned:
\begin{itemize}
\item predicted class SUPPORTED: use \texttt{attention\_sup},
\item predicted class REFUTED: use \texttt{attention\_ref},
\item if no target class is passed: use $\max(\texttt{attention\_sup},\texttt{attention\_ref})$.
\end{itemize}
Default rationale budget: \texttt{rationale\_top\_n = 10}.

\subsection{Faithfulness checks implemented}
Two complementary checks are available:
\begin{enumerate}
\item \textbf{Leave-one-out (LOO):} remove one rationale edge at a time and measure confidence drop for predicted class.
\item \textbf{Grouped removal:} remove top-$k$ rationale edges together and measure confidence drop.
\end{enumerate}

Current boolean criterion is positive confidence drop:
\[
\texttt{is\_faithful} = [\Delta\text{confidence} > 0].
\]
M3 runs faithfulness only on the first \texttt{faithfulness\_subset\_size} claims (default $3$) to keep runtime bounded.

\subsection{What M3 logs for interpretability}
Per claim, M3 logs:
\begin{itemize}
\item \texttt{top\_rationale\_edges} with attention, PV relevance, and combined score,
\item \texttt{faithfulness\_loo} with per-edge confidence drops,
\item \texttt{salience\_results} (group removal) with aggregate drop.
\end{itemize}
Run-level summary adds counts such as:
\texttt{num\_faithful\_claims} and \texttt{num\_salience\_checked}.
